% concatenated from main.tex
\documentclass[a4paper]{amsart}
\usepackage[utf8]{inputenc}
\usepackage[english]{babel}
\usepackage{fancyhdr}
%\pagestyle{fancy}
\usepackage{amsmath}
\usepackage{amsfonts,amstext}
\usepackage{amssymb,amsthm,amscd,amsxtra,wasysym,graphicx}
%\usepackage[OT1]{fontenc}
\usepackage{mathrsfs,esint,comment}
%\usepackage{tikz-cd}
\usepackage{enumitem,mathtools,dsfont}
       \usepackage{palatino,mathpazo}
  \usepackage{palatino}
\usepackage{charter}
\usepackage{xcolor}
\def\blue{\color{blue}}
\def\orange{\color{orange}}
\def\red{\color{red}}
\def\green{\color{green}}  

\def\Bibtex{{\rm B\kern-.05em{\sc i\kern-.025em b}\kern-0.08em T\kern-.1667em\lower.7ex\hbox{E}\kern-.125emX}}
\hfuzz1pc
    % \usepackage[breaklinks=true,bookmarks=false,pagebackref]{hyperref}
\usepackage{hyperref}
\hypersetup{
	unicode=false,          
	pdftoolbar=true,       
	pdfmenubar=true,       
	pdffitwindow=false,     
	pdfstartview={FitH}, 
	pdftitle={Monge-Amp\`ere equations},   
	pdfauthor={Quang-Tuan Dang},   
	colorlinks=true,  
	linkcolor=purple,         
	citecolor=blue,        
	filecolor=green,      
	urlcolor=blue}          

% \textwidth=125mm
% \textheight=185mm
% \parindent=8mm
% \evensidemargin=0pt
% \oddsidemargin=0pt
% \frenchspacing
\usepackage{geometry}
\usepackage{fullpage}
\usepackage{indentfirst}

\usepackage{colonequals}
\usepackage{enumitem}
% \usepackage{refcheck}
\renewcommand{\baselinestretch}{1.05}
\numberwithin{equation}{section}

%\usepackage{showkeys}


\newtheorem{thmA}{Theorem}
\renewcommand{\thethmA}{\Alph{thmA}}
\newtheorem{corA}[thmA]{Corollary}
%\renewcommand{\thecorA}{\Alph{corA}}

% theorems with special labels
\newtheorem{thmspec}{\relax}

\newtheorem{theorem}{Theorem}[section]
\newtheorem{proposition}[theorem]{Proposition}
\newtheorem{corollary}[theorem]{Corollary}
\newtheorem{lemma}[theorem]{Lemma}
\theoremstyle{definition}
\newtheorem{definition}[theorem]{Definition}
\newtheorem{remark}[theorem]{Remark}
\newtheorem{remarks}[theorem]{Remarks}
\newtheorem{example}[theorem]{Example}
\newtheorem{exercise}[theorem]{Exercise}
\newtheorem{question}[theorem]{Question}
\newtheorem{conjecture}[theorem]{Conjecture}
\newtheorem{problem}[theorem]{Problem}

\newcommand{\cali}[1]{\mathscr{#1}}

\newcommand{\GL}{{\rm GL}}
\newcommand{\Sp}{{\rm Sp}}
\renewcommand{\sp}{{\rm \mathfrak{sp}}}
\newcommand{\SO}{{\rm SO}}
\newcommand{\SU}{{\rm SU}}
%\newcommand{\U}{{\rm U}}
\newcommand{\End}{{\rm End}}
\newcommand{\Vol}{{\rm Vol}}
\newcommand{\Hol}{{\rm Hol}}
\newcommand{\SL}{{\rm SL}}
\newcommand{\MA}{{\rm MA}}
\newcommand{\NMA}{{\rm NMA}}
\renewcommand{\sl}{{\rm\mathfrak{sl}}}
\newcommand{\Aut}{{\rm Aut}}
\newcommand{\lov}{{\rm lov}}
\newcommand{\supp}{{\rm Supp}}
\newcommand{\wi}{\widetilde}
\newcommand{\maxep}{\mathop{\mathrm{\rm  max_\epsilon}}\nolimits}
\newcommand{\maxunit}{\rm  max_1}
\newcommand{\const}{\mathop{\mathrm{const}}\nolimits}
\newcommand{\ev}{{\rm ev}}
\newcommand{\ov}{\widehat}
\newcommand{\Int}{{\rm Int}}

\newcommand{\diff}{\text{\normalfont d}}

\newcommand{\Leb}{\mathop{\mathrm{Leb}}\nolimits}
\newcommand{\Lip}{\mathop{\mathrm{Lip}}\nolimits}
\newcommand{\LLip}{{\mathop{\mathrm{\widetilde{Lip}}}\nolimits}}

\newcommand{\dist}{\mathop{\mathrm{dist}}\nolimits}
\renewcommand{\Re}{\mathop{\mathrm{Re}}\nolimits}
\renewcommand{\Im}{\mathop{\mathrm{Im}}\nolimits}
\newcommand{\vol}{\mathop{\mathrm{vol}}}

\newcommand{\eq}{\mathrm{eq} }
%\newcommand{\ddc}{{dd^c}}
\newcommand{\ddc}{dd^c}
\newcommand{\dc}{d^c}
%\newcommand{\d}{\text{\normalfont d}}
%\newcommand{\Ta}{\text{\normalfont T}}
\newcommand{\Fix}{\text{\normalfont Fix}}
\newcommand{\Reg}{\text{\normalfont Reg}}
\newcommand{\Sing}{\text{\normalfont Sing}}






\newcommand{\mult}{{\rm mult}}
\newcommand{\Ssupp}{{\rm supp^*}}
\newcommand{\loc}{{loc}}
%\newcommand{\ddc}{dd^c}
%\newcommand{\dc}{d^c}
\newcommand{\diag}{{\rm diag}}
\newcommand{\prim}{{\rm prim}}
\newcommand{\otextbf}{{\rm \mathbf{0}}}


\frenchspacing

\textwidth=13.5cm
\textheight=23cm
\parindent=16pt
\oddsidemargin=1.1cm
\evensidemargin=1.1cm
\topmargin=-0.5cm

\def\B{\mathbb B}
\def\und{\underset}
\def\O{\Omega}
\def\wed{\wedge}
\def\n{\noindent}
\def\C^n{\mathbb C^n}
\def\var{\varepsilon}
\def\te{\text}
\def\be{\beta}
\def\de{\delta}
\def\lam{\lambda}
\def\Cal{\mathcal}
\def\al{\alpha}
\def\lg{\longrightarrow}
\def\cen{\centerline}
\def\R{\mathbb{R}}
\frenchspacing
%\DeclareMathOperator{\loc}{loc}
\numberwithin{equation}{section}

\def\cB{\mathcal B}
\def\bC{\mathbb C}
\def\cE{\mathcal E}
\def\cF{\mathcal F}
\def\cO{\mathcal O}
\def\cH{\mathcal H}
\def\cI{\mathcal I}
\def\cJ{\mathcal J}
\def\n{\noindent}

\newcommand{\capK}{\text{Cap}}
%\newcommand{\span}{\rm span}

\newcommand{\PSH}{{\rm PSH}}
%\newcommand{\Z}{\mathbb{Z}}
\newcommand{\FS}{{\rm FS}}
%\newcommand{\B}{\mathbb{B}}
%\newcommand{\C}{\mathbb{C}}
\newcommand{\D}{\mathbb{D}}
\newcommand{\N}{\mathbb{N}}
\newcommand{\Z}{\mathbb{Z}}
%\newcommand{\R}{\mathbb{R}}
\renewcommand\P{\mathbb{P}}

\renewcommand{\S}{\mathbb{S}}

\newcommand{\e}{\textbf{\textit{e}}}



%\newcommand{\K}{{\cal K}}
\newcommand{\E}{{\mathbb{E}}}
\renewcommand{\H}{\mathbb{ H}}

\newcommand{\swedge}{\widehat\wedge}
\newcommand{\transposee}[1]{{\vphantom{#1}}^{\mathit t}{#1}}
\newcommand{\spin}{$\text{spin}^c$ }
\newcommand{\norm}[1]{\lVert#1\rVert}
\newcommand{\abs}[1]{\lvert#1\rvert}

%------------------------------Boldsymbol-------------------------------------
\newcommand{\boldsym}[1]{\boldsymbol{#1}}
\newcommand\bb{\boldsym{b}}
\newcommand\bE{\boldsym{E}}
\newcommand\bg{\boldsym{g}}
\newcommand\bI{\boldsym{I}}
\newcommand\bJ{\boldsym{J}}
\newcommand\bk{\boldsym{k}}
\newcommand\bm{\boldsym{m}}
\newcommand\br{\boldsym{r}}
\newcommand\bt{\boldsym{t}}


\newcommand*{\doi}[1]{\href{http://dx.doi.org/#1}{doi: #1}}
\setcounter{tocdepth}{1}
%\setcounter{secnumdepth}{3}
%opening

\usepackage{microtype}
\begin{document}

\title{Singularities vs non-pluripolar Monge--Amp\`{e}re masses}
\author{Quang-Tuan Dang, Hoang-Son Do, Hoang Hiep Pham}


\address{The Abdus Salam International Centre for Theoretical Physics (ICTP), Str. Costiera, 11, 34151 Trieste, TS, Italy}
\email{qdang@ictp.it}
\address{Vietnam Academy of Science and Technology, Institute of Mathematics, 18 Hoang Quoc Viet Road, Cau Giay, Hanoi, Vietnam}
\email{dhson@math.ac.vn $\&$ phhiep@math.ac.vn}
% \address{Vietnam Academy of Science and Technology, Institute of Mathematics, 18 Hoang Quoc Viet Road, Cau Giay, Hanoi, Vietnam}
% \email{phhiep@math.ac.vn}
\date{\today}
\subjclass[2020]{32W20, 32U05, 32Q15}
\keywords{Complex Monge-Amp\`ere operators, Singularity type, Capacity}


%\renewcommand{\thefootnote}{}
%\footnote{2010 \emph{Mathematics Subject Classification}: 32C25, 32C30, 32U25, 32W20.}
%\vskip0.1cm

%\footnote{\emph{Key words and phrases}: a pure $k$-dimentional analytic subsets, plurisubharmonic functions on analytic subsets, Lelong number, the log canonical thresholds.}

%\renewcommand{\thefootnote}{\arabic{footnote}}
%\renewcommand{\thefootnote}{\arabic{footnote}}
%\setcounter{footnote}{0}


\begin{abstract}
 The aim of this paper is to compare singularities of closed positive currents whose non-pluripolar complex Monge--Ampère masses equal. We also provide a short alternative proof for the monotonicity of non-pluripolar complex Monge--Ampère masses, generalizing results of Witt-Nystr\"om, Darvas--Di Nezza--Lu, Lu--Nguy\^en and Vu. 
 %introduce and study the notion of the complex Monge-Amp\`{e}r measure of an arbitrary plurisubharmonic function.
\end{abstract}

\maketitle
% \tableofcontents
\section{Introduction}


One of the central topics of interest in pluripotential theory is the domain of the definition of the complex Monge--Amp\`ere operator. 
% Since the 1970s, Bedford and Taylor constructed the Monge--Amp\`{e}re operator on the class of locally bounded plurisubharmonic functions and investigated the existence of solutions to the Dirichlet problem on bounded strictly pseudoconvex domains (see \cite{bedford1976dirichlet}). Subsequently, in their seminal work (see \cite{bedford1982new}), they established profound properties of plurisubharmonic functions on open subsets of $\mathbb{C}^n$. The Monge–Ampère operator, as defined by Bedford and Taylor (\cite{bedford1976dirichlet}, \cite{bedford1982new}), applies to locally bounded plurisubharmonic functions. 
Building on the seminal work of Bedford and Taylor \cite{bedford1976dirichlet,bedford1982new}, the Monge--Amp\`{e}re operator is well-defined on the class of locally bounded plurisubharmonic functions.
This was later extended by Demailly~\cite{demailly1993monge} to plurisubharmonic functions with isolated or compactly supported singularities. Furthermore, Cegrell~\cite{cegrell1998pluricomplex,cegrell2004general} introduced the largest class $\mathcal{E}$ of plurisubharmonic functions, on which the Monge--Amp\`ere operator is well-defined and continuous along decreasing sequences.

% Non-pluripolar Monge--Ampère products of positive closed $(1,1)$-currents on a compact K\"{a}hler manifold was first introduced by Guedj--Zeriahi \cite{guedj2007weighted} and 

 Guedj and Zeriahi~\cite{guedj2007weighted} extended the Bedford–Taylor theory~\cite{bedford1987fine} and Cegrell theorey~\cite{cegrell1998pluricomplex} (originally developed in a local setting) to the geometric setting of compact K\"ahler manifolds, where they defined the so-called non-pluripolar Monge--Amp\`ere operator for unbounded potentials. Their approach was later adapted to the setting of big cohomology classes by Boucksom, Eyssidieux, Guedj, and Zeriahi \cite{boucksom2010monge}. In the latter paper, the authors posed a conjecture regarding the monotonicity of Monge--Ampère masses, which was later resolved by D. Witt-Nystr\"om  \cite{wittnystrom19-monotonicity}. Alternative proofs were subsequently provided in \cite{vu2021relative,lu2022hessian,lu2021capacities}.
In this paper, we generalize this result by showing that the complex Monge--Ampère mass of potentials decreases as its singularity type increases {\em in capacity}; cf. Definition~\ref{def: sing-capa}. 
%We observed that a version of this conjecture in the local setting was Theorem 4.1 in \cite{Ahag09-MA-pluripolar}. In this note, we will use ideas in the proof of Theorem 4.1 in \cite{Ahag09-MA-pluripolar} to give a direct proof of a following result about the monotonicity of non-pluripolar products which improve Theorem 1.2 in \cite{wittnystrom19-monotonicity} and Theorem 1.1 in \cite{darvas2023relative}.

To state our results, let $(X, \omega)$ be a compact K\"{a}hler manifold of dimension $n$ and let $\theta$ be a smooth closed real(1,1) form. 
% We say that $u\in L^1(X)$ is $\theta$-plurisubharmonic ($\theta$-psh)  if, locally, $u$ can be written as the sum of a plurisubharmonic and a smooth function, and $\theta_u:=\theta+\ddc u\geq 0$ in the sense of currents. 
We let $\PSH(X, \theta)$ denote the set of $\theta$-psh functions on $X$. Recall that
the cohomology class $\{\theta\}$ is {\em pseudoeffective} if $\PSH(X,\theta)\neq \varnothing$. 
 The cohomology class $\{\theta\}$ is is {\em big} if $\PSH(X,\theta-\varepsilon_0\omega)\neq \varnothing$ for some $\varepsilon_0>0$. 

 If $u$ and $v$ are two $\theta$-psh functions on X, then $u$ is said to be {\em less (resp. more) singular} than $v$ if $u  \geq\, v+O(1)$ (resp. $u  \leq\, v+O(1)$). 
 We say that $u$ has {\em the same singularities} as $v$ if $u$ is less singular than $v$, and $v$ is less singular than $u$. 
% This defines an equivalence relation on $\PSH(X,\theta)$, whose equivalence classes are the singularity type $[u]$.

% Let $T_j=\theta_j+\dc u_j$, for $j\in\{1,\ldots,p\}$, $1\leq p\leq n$, be closed positive currents, where $\theta_j$ are closed smooth real (1,1)-forms. Boucksom--Eyssidieux--Guedj--Zeriahi \cite{boucksom2010monge} defined the {\em non-pluripolar product} of these currents: 
% $$\theta_{1,u_1}\wedge\ldots\wedge \theta_{p,u_p}=\langle T_1\wedge\ldots \wedge T_p\rangle$$
% which is shown to be a well-defined closed positive $(p,p)$-current on $X$ putting no mass on pluripolar sets. 
For a $\theta$-psh function $u$, one defines its {\em non-pluripolar Monge--Amp\`ere measure} by $\theta_u^n:=\theta_u\wedge\ldots\wedge\theta_u$. Fix $\phi\in \PSH(X,\theta)$.
% We say that $u$ has {full mass relative} to $\phi$ if $u$ is more singular than $\phi$ and $\int_X\theta^n_u=\int_X\theta_\phi^n$. 
 % Let $\PSH(X,\theta,\phi)$ denote the set of $\theta$-psh functions that are more singular than $\phi$.
Let
$\mathcal{E}(X,\theta,\phi)$ denote the set of  $\theta$-ph functions $u$ with full mass relative to $\phi$, i.e., $u$ is more singular than $\phi$ and $\int_X\theta^n_u=\int_X\theta_\phi^n$. When $\phi=V_\theta$ is the "the least singular" element of $\PSH(X,\theta)$, we simply denote by $\mathcal{E}(X,\theta)$.

% Witt Nystr\"om showed in \cite{wittnystrom19-monotonicity} the dependence of the Monge--Ampère mass on the singularity type. Specifically, if $\varphi,\psi\in\PSH(X,\theta)$ have the same singularity type, i.e, $|\varphi-\psi|\leq O(1)$, then their Monge--Ampère masses are equal: $\int_X\theta_\varphi^n=\int_X\theta_\psi^n$. 
%  A natural question arises: conversely, if $\int_X\theta_\varphi^n=\int_X\theta_\psi^n$, what can one say about the relationship between the singularity types of $\varphi$ and $\psi$? Our main result is the following.

Extending the result of Guedj--Zeriahi~\cite{guedj2007weighted} and Dinew~\cite{dinew2009uniqueness} for the K\"ahler case, Darvas--Di Nezza--Lu~\cite{darvas2018monotonicity,darvas2021log} 
 showed that the non-pluripolar Monge--Ampère measure uniquely determines the potential within a relative full mass class $\mathcal{E}(X,\theta,\phi)$. 
% gave a complete characterization of the range of the complex Monge--Ampère operator on the relative full mass class $\mathcal{E}(X,\theta,\phi)$.
They also provided a characterization of membership in the class $\mathcal{E}(X,\theta,\phi)$, solely in terms of singularity type. Our first theorem below gives another characterization.
\begin{theorem}\label{thm: thm1}
	%Let $\varphi,\psi\in\PSH(X,\theta)$ with positive masses. If $\int_X\theta_\varphi^n=\int_X\theta_\psi^n>0$, 
Let $(X, \omega)$ be a compact K\"{a}hler manifold and $\theta$ a smooth closed real (1, 1)-form on $X$ whose cohomology class is big.
Let $\phi\in\PSH(X,\theta)$. If $\varphi,\psi\in\mathcal{E}(X,\theta,\phi)$, then for each $c>0$, there exists a function $ h\in\mathcal{E}(X,\omega)$ such that $$|\varphi-\psi|\leq -ch.$$
Conversely, if $\varphi,\psi\in\PSH(X,\theta)$ such that $|\varphi-\psi|\leq -ch$ for some $c>0$ and $h\in\mathcal{E}(X,\omega)$, then both $\varphi$ and $\psi$ belong to $\mathcal{E}(X,\theta,\max(\varphi,\psi))$.

In particular, $\varphi\in\mathcal{E}(X,\theta,\phi)$ if and only if $\varphi$ is more singular than $\phi$ and $\varphi\geq \phi+ch $ for some $c>0$ and $h\in\mathcal{E}(X,\omega)$.
% $$\int_X\theta_\varphi^n=\int_X\theta_\psi^n.$$
% Furthermore, if $\theta_\varphi^n=\theta_\psi^n$ then $u-v$ is constant. 
\end{theorem}  
% The last statement shows that the non-pluripolar Monge--Ampère measure uniquely determines the potential within a relative full mass class. This also establishes the uniqueness result of Dinew~\cite{dinew2009uniqueness} for the K\"ahler case, in which his method can be adapted to more singular settings ~\cite{boucksom2010monge,darvas2018monotonicity}. 
% We are, moreover, able to avoid the mass concentration technique, which 
The first statement of this theorem is a direct consequence of a result of the first author~\cite{dang2021continuity} (see Proposition~\ref{prop: homega}), whose proof relies on the envelope technique, developed in the series of recent works~\cite{darvas2018singularity,darvas2018monotonicity,darvas2020metric}. 
% the plurifine locality property~\cite{bedford1987fine} and the monotonicity of non-pluripolar Monge--Ampère masses proved by Witt-Nystr\"om \cite{wittnystrom19-monotonicity}; see also \cite{vu2021relative,lu2022hessian,lu2021capacities} for different proofs. 
For the second statement,
 inspired by the idea in \cite{Ahag09-MA-pluripolar}, we establish a slight generalization of the monotonicity of non-pluripolar Monge--Ampère masses, as shown in \cite{wittnystrom19-monotonicity} (see also \cite{vu2021relative,lu2022hessian,lu2021capacities} for different proofs), thereby providing an alternative proof of Witt-Nystr\"om's result.
\begin{theorem}\label{mainthm}
Let $(X, \omega)$ be a compact K\"{a}hler manifold of dimension $n$. Let $\theta_j$, $j\in\{1,\ldots,n\}$, be smooth closed real $(1,1)$-forms on $X$ whose cohomology classes are pseudoeffective. Assume $\varphi_j,\psi_j \in \PSH (X,\theta_j )$. If there exist $h \in \mathcal E (X, \omega )$ and $c\geq 0$ such that 
$$\lim\limits_{ k\to+\infty }k^n \capK_{\omega} ( \{ \varphi_j \ < \psi_j+ c h - k \} ) = 0,$$
for all $1\leq j\leq n$, then 
$$\int_X  \theta_{ 1, \varphi_1 }\wedge ...\wedge \theta_{ n, \varphi_n }  \geq \int_X  \theta_{ 1, \psi_1 }\wedge ...\wedge \theta_{ n, \psi_n }  .$$
\end{theorem}
Here, the notation of Monge--Ampère capacity $\capK_\omega$ was introduced by S. Ko\l odziej \cite{kolodziej2003monge}: for any Borel set $E\subset X$, one defines
$$\capK_{ \omega }(E):=\sup\Big\{\int_X\omega_u^n: u\in\PSH(X,\omega),\, 0\leq u\leq 1 \Big\}.$$ 


%The class $\mathcal E (X, \omega )$ of $\omega$-plurisubharmonic functions with finite weighted Monge-Am\`{e}pere energy was introduced and studied by Vincent Guedj and Ahmed Zeriahi \cite{guedj2007weighted}.


% The above main theorem directly implies a comparison principle as follows:
% \begin{corollary}
% Let $(X,\omega)$ be a compact K\"{a}hler manifold of dimension $n$. Let $\theta$ be a smooth closed real (1, 1)-form on $X$ whose cohomology class is pseudoeffective. Let $\varphi , \psi \in \PSH (X,\theta )$. If there exist $h\in \mathcal E (X, \omega )$ and $c\geq 0$ such that 
% $$\lim_{ k\to+\infty } k^n \capK_{ \omega } ( \{ \varphi < \psi + c h - k \} ) = 0$$
% then 
% $$\int_{\{ \varphi < \psi\} }  \theta_{ \psi }^n  \leq \int_{ \{ \varphi < \psi \}}  \theta_{ \varphi }^n  .$$
% \end{corollary}




\subsection*{Acknowledgement} We are grateful to Hoang-Chinh Lu for his  comments on
a preliminary draft.  
This work was done while the authors were visiting Vietnam Institute for Advanced Study in Mathematics (VIASM), and we would like to thank VIASM for its hospitality.

\subsection*{Ethics declarations} The authors declare no conflict of interest.
\section{Preliminaries}
In the whole article, we let $X$ denote a compact K\"ahler manifold of dimension $n$, equipped with a K\"ahler form $\omega$, and $\theta$ a smooth closed real $(1,1)$-form on $X$. Without loss of generality, we may assume that $\theta\leq \omega$.
\subsection{Non-pluripolar products}
Recall that a function $u:X\to\mathbb{R}\cup\{-\infty\}$ is quasi-plurisubharmonic
(quasi-psh) if, locally, $u$ can be written as the sum of a plurisubharmonic function and a smooth function. 
We say that a function $u$ is $\theta$-plurisubharmonic ($\theta$-psh) if it is quasi-psh and $\theta_u:=\theta+\ddc u\geq 0$ in the sense of currents.  We let $\PSH(X, \theta)$ denote the set of $\theta$-psh functions on $X$. Recall that
the cohomology class $\{\theta\}$ is {\em pseudoeffective} if $\PSH(X,\theta)\neq \varnothing$. 
 The cohomology class $\{\theta\}$ is is {\em big} if $\PSH(X,\theta-\varepsilon_0\omega)\neq \varnothing$ for some $\varepsilon_0>0$.

 
Let $x_0\in X$. Fixing a holomorphic chart $x_0\in V_{x_0}\subset X$,  the {\em Lelong number} $\nu(\varphi,x_0)$ of a quasi-psh function $\varphi$ at  $x_0\in X$ is defined as follows:
		\begin{align*}
		\nu(\varphi,x_0):=\sup\{\gamma\geq 0: \varphi(z)\leq \gamma\log\|z-x_0\|+O(1), \; \text{on}\; V_{x_0}\}.  
		\end{align*}
 Remark that this definition does not depend on the choice of local holomorphic charts.

   % For $\varphi\in\PSH(X,\theta)$, one defines {\em the multiplier ideal sheaf} $\mathcal{I}(t\varphi,x_0)$, $t\geq 0$, to be the sheaf of germs of holomorphic functions $f$ such that $|f|^2e^{-t\varphi}$ is locally integrable with respect to the Lebesgue measure in some local coordinates near $x_0$.

\smallskip

Assume $u,v\in \PSH(X,\theta)$. We say that $u$ is {\it less (resp. more) singular} than $v$, and denote by $u\succeq v$ (resp. $u \preceq v$), if there exists a constant $C$ such that $u+C\geq v$ (resp. $u\leq v+C$) on $X$. We say that $u$, $v$ have the {\em same singularity type}, and denote by $u\simeq v$ if $u \preceq v$ and $u\succeq v$.
A $\theta$-psh function $u$ is said to have {\em minimal singularities} if it is less singular than any other $\theta$-psh function.
We introduce the extremal function $V_\theta$ defined by
\[ V_\theta(x):=\sup\{\varphi(x)\in\PSH(X,\theta): \varphi\leq 0 \}.\]
One can see that $V_\theta$ is a $\theta$-psh function with minimal singularities. 
In particular, if $\theta$ is semi-positive, then $V_\theta=0$.

Let $\theta_1,\ldots,\theta_p$ be smooth closed real (1,1)-forms for $1\leq p\le n$. 
Let $u_j$ be $\theta_j$-psh functions for $j\in\{1,\ldots,p\}$ and put $\theta_{j,u_j}:=\theta_j +\ddc u_j$. We recall how to define the non-pluripolar product $\theta_{u_1} \wedge \cdots \wedge \theta_{u_n}$. We write locally $\theta_{j,u_j}= \ddc v_j$, where $v_j$ is psh. By \cite{bedford1987fine, boucksom2010monge}, one knows that the sequence of positive currents 
$$\mathbf{1}_{\cap_{j=1}^p \{v_j>-k\}}\ddc \max\{v_1, -k\} \wedge \cdots \wedge \ddc \max\{v_p, -k\}$$ 
is increasing in $k \in \N$ and converges to a closed positive current, which is independent of the choice of local potentials $v_j$'s. Thus, we obtain a well-defined closed positive current on $X$ which is called the \emph{non-pluripolar product} $\theta_{1,u_1} \wedge \cdots \wedge \theta_{p,u_p}$ of $\theta_{1,u_1}, \ldots, \theta_{p,u_p}$.  
%If $p=n$, the resulting positive $(n, n)$-current is a Borel measure putting no mass on pluripolar sets. 
For any $u\in\PSH(X,\theta)$, the {\em non-pluripolar complex Monge--Amp\`ere measure} of $u$ is given by \[\theta_u^n:= (\theta+\ddc u)^n.\]
Given a potential $\phi\in\PSH(X,\theta)$, we let $\PSH(X,\theta,\phi)$ denote the set of $\theta$-psh functions $u$ such that $u\leq\phi$. We denote by $\mathcal{E}(X,\theta,\phi)$ the set of functions $u\in\PSH(X,\theta,\phi)$ with {\em full Monge--Amp\`ere mass relative} to $\phi$, i.e.,
$\int_X\theta_u^n=\int_X\theta_\phi^n$.
If $\phi=V_\theta$ we simply denote by $\mathcal{E}(X,\theta)$. The volume of $\theta$ is defined by $\Vol(\theta):=\int_X\theta_{V_\theta}^n.$
We remark that by~\cite[Theorem 4.7]{boucksom2002volume}, the class $\{\theta\}$ is big if and only if $\Vol(\theta)>0$


% For every Borel set $E$ in $X$, recall that the capacity of $E$ is given by 
% $$\capK_{ \omega }(E):=\sup\Big\{\int_X\omega_u^n: u\in\PSH(X,\omega),\, 0\leq u\leq 1 \Big\}.$$

	We recall here the plurifine locality of the non-pluripolar Monge--Amp\`ere product (see~\cite[Corollary 4.3]{bedford1987fine} or \cite[Section 1.2]{boucksom2010monge}) for later use. This topology is the coarsest such that all quasi-psh functions with values in $\mathbb{R}$ are continuous.
\begin{lemma}\label{lem: plurifine}
	Assume that $\varphi$, $\psi$ are $\theta$-psh functions such that $\varphi=\psi$ on an open set $U$ in the plurifine topology. Then 
	\begin{equation*}
	\mathbf{1}_U\theta_\varphi^n=\mathbf{1}_U\theta_\psi^n.
	\end{equation*}
\end{lemma}
For practice, we stress that sets of the form
$\{u < v\}$, where $u$ and $v$ are quasi-psh functions, are open in the plurifine topology. Lemma \ref{lem: plurifine} will be referred to as the {\em plurifine locality property}.


% We recall the following classical inequality.
\begin{lemma}\label{lem: maxprin}
	Let $\varphi,\psi\in\PSH(X,\theta)$. Then
	\[\theta_{\max(\varphi,\psi)}^n\geq \mathbf{1}_{\{\psi\leq\varphi \}}\theta_\varphi^n+\mathbf{1}_{\{\varphi<\psi\}}\theta_{\psi}^n. \]
    In particular, if $\varphi\leq\psi$ then $\mathbf{1}_{\{\varphi=\psi\}}\theta^n_\varphi\leq \mathbf{1}_{\{\varphi=\psi\}}\theta^n_\psi$.
\end{lemma}
% \begin{proof} We refer the readers to \cite[Lemma 2.5]{darvas2023relative}.
	%For the reader's convenience, we briefly give proof here. 
%	If $\varphi$ and $\psi$ are locally bounded, the result follows from the maximum principle; see e.g.,~\cite[Theorem 3.23]{guedj2017degenerate}. For general case, set $\varphi^t:=\max(\varphi,V_\theta-t)$, $\psi^t:=\max(\psi,V_\theta-t)$. Then $\varphi^t$ and $\psi^t$ are locally bounded on $\Omega$, it follows that
%	\[\mathbf{1}_{\{\varphi>V_\theta-t\}\cap\{\psi>V_\theta-t \}\cap \{\varphi^t=\psi^t\}}\theta^n_{\varphi^t}\leq \mathbf{1}_{\{\varphi>V_\theta-t\}\cap\{\psi>V_\theta-t \}\cap \{\varphi^t=\psi^t\}}\theta^n_{\psi^t}, \]
	%multiplying with $\mathbf{1}_{\{\f>V_\theta-t\}\cap\{\psi>V_\theta-t \}}$ both side, and 
%	using plurifine locality. Letting $t\to+\infty$, the inequality follows easily.
	%We refer the reader to~\cite{darvas2020metric}.
% \end{proof}

\subsection{Quasi-envelopes}
Given a measurable function $f:X\to\mathbb{R}$, we define the {\em $\theta$-psh envelope} of $h$ by
\begin{equation*}
P_\theta(f):=(\sup\{u\in\PSH(X,\theta): u\leq f\;\text{on}\, X \})^*,
\end{equation*} where the star means we take the upper semi-continuous regularization. %We next define 
%\[ P_{\theta}(g,h):=(\sup\{u\in\PSH(X,\theta): u\leq \min (g,h)\;\text{on}\, X \})^*.\] 

Given a $\theta$-psh function $\phi$, 
Ross and Witt Nystr\"om~\cite{ross2014analytic} introduced the "rooftop envelope" as follows
\[P_\theta[\phi](f)= \left( \lim_{C\to +\infty}P_\theta(\min(\phi+C,f))\right)^*. \]
When $f=0$ we simply write $P_\theta[\phi]$. 
\begin{definition}
A function $\phi\in\PSH(X,\theta)$ is called a {\em model potential} if $\int_X\theta_\phi^n>0$ and $\phi=P_\theta[\phi]$.
\end{definition}
We recall the first technical result about quasi-envelopes. 
\begin{theorem} \label{thm: envelope}
	Assume that $f$ is a quasi-continuous function and $P_\theta(f)\in\PSH(X,\theta)$. Then, $P_\theta(f)\leq f$ outside a pluripolar set, and $\theta_{P_\theta(f)}^n$ is concentrated
	on the contact set $\{P_\theta(f)=f \}$.
\end{theorem}
\begin{proof}
 See \cite[Theorem 2.2]{darvas2023relative}.
\end{proof}

\begin{proposition}\label{prop: homega}
Let $\{\theta\}$ be a big cohomology class and $\phi\in\PSH(X,\theta)$. Let $\varphi\in\mathcal{E}(X,\theta,\phi)$ and $\psi\in \PSH(X,\theta,\phi)$. Then for any $b>0$, $P_\omega(b\varphi-b\psi)\in\mathcal{E}(X,\omega)$.
\end{proposition}
\begin{proof}
The proof was given in~\cite[Proposition 2.10]{dang2021continuity} in the special case when $\phi=V_\theta=\psi$. The proof is based on Theorem \ref{thm: envelope}, the plurifine locality (Lemma \ref{lem: plurifine}), and the monotonicity of Monge--Amp\`ere masses (Theorem~\ref{mainthm}).
The general case, where $\phi$ is an arbitrarily $\theta$-psh function with positive mass, follows by the same argument.
\end{proof}
As a consequence, we confirm the question posed by Berman--Boucksom--Eyssidieux--Guedj--Zeriahi~\cite[Question 36]{dinew2016open}, which was previously resolved by Darvas--Di Nezza--Lu~\cite{darvas2018singularity}.
\begin{corollary}
Assume that $\{\theta\}$ is a big and nef cohomology class and $\varphi\in\mathcal{E}(X,\theta)$. Then, $\varphi$ has zero Lelong number at all points. 
 % Moreover, $\mathcal{I}(\varphi,x)=\mathcal{I}(V_\theta,x)$.
\end{corollary} 
\begin{proof}
Let $\varphi\in \mathcal{E}(X,\theta)$. 
We apply Proposition~\ref{prop: homega} with $b=1$, $\phi=\psi=V_\theta$ to obtain $h\colonequals P_\omega(\varphi-V_\theta)\in\mathcal{E}(X,\omega)$. By Theorem~\ref{thm: envelope}, we have $V_\theta+h\leq\varphi$ almost everywhere, and hence everywhere on $X$. It follows that $\nu(\varphi,x)=\nu(V_\theta,x)$ for all $x\in X$. From~\cite[Propositions 3.2, 3.6]{boucksom2004divisorial}, we have $\nu(V_\theta,x)=0$ for all $x\in X$, which completes the proof. 
% Since $\varphi\geq V_\theta+h$
% $\varphi\in\mathcal{E}(X,\omega)$, and the conclusion follows immediately from ~\cite[Corollary 1.8]{guedj2007weighted}.
\end{proof}
We also point out a slight generalization of Proposition~\ref{prop: homega}.
\begin{proposition}\label{prop: hOmega}
	\label{prop_imp}
	Let $\{\theta\}$ be a big cohomology class. Let $\varphi,\psi\in\PSH(X,\theta)$ be such that $P_\theta[\varphi]$ is less singular than $\psi$. Then for any $b>0$, $P_\omega(b\varphi-b\psi)\in\mathcal{E}(X,\omega)$.  
\end{proposition}
\begin{proof}
	We refer the reader to \cite[Proposition 2.4]{dang2023continuity}. 
    % The proof is based on Theorem \ref{thm: envelope}, the plurifine locality (Lemma \ref{lem: plurifine}), and the monotonicity of Monge--Amper\`re masses. 
\end{proof}
% We end this section with the following definition.
\begin{definition}\label{def: sing-capa}
Assume $\varphi,\psi\in\PSH(X,\theta)$. We say that $\varphi$ is {\em less singular in capacity} than $\psi$, denoted by $\varphi\succeq_C \psi$, if there exist $h\in\mathcal{E}(X,\omega)$ and $c\geq 0$ such that
\begin{equation}\label{eq: def-sing-capa} \lim_{k\to+\infty}k^n\capK_{\omega}(\{\varphi<\psi+ch-k\})=0.\end{equation}
We say $\varphi$, $\psi$ have the {\em same singularity
type in capacity} if $\varphi\succeq_C \psi$ and $\psi\succeq_C \varphi$.
\end{definition}
 
\section{Proof of main theorems}
% \subsection{Monotonicity of non-pluripolar products}
% From now on, $\theta$ always denotes a smooth closed (1,1)-form whose cohomology class is pseudoeffective.

First, we prove the following proposition, which states that potentials with the same singularity type also have the same global mass. 
This is a conjecture of Boucksom--Eyssidieux--Guedj-Zeriahi~\cite[page 217]{boucksom2010monge}, which was initially proved by Witt--Nystr\"om~\cite{wittnystrom19-monotonicity}.
% This result generalizes that of Witt--Nystr\"om~\cite{wittnystrom19-monotonicity}, who proved athe conjecture of Boucksom--Eyssidieux--Guedj-Zeriahi~\cite[page 217]{boucksom2010monge}. 


\begin{proposition}\label{prop1}
Let $\varphi,\psi\in\PSH(X,\theta)$. If $\varphi$ and $\psi$ have the same singularity type, then $\int_X\theta_\varphi^n=\int_X\theta_\psi^n$.  
\end{proposition}
We would like to mention that several alternative proofs of this conjecture were later given by Lu--Nguy\^en~\cite{lu2022hessian} using the monotonicity of the energy functional, Lu~\cite{lu2021capacities} using a standard approximation process, and Vu~\cite{vu2021relative} where relative non-pluripolar products of
currents are constructed. We provide below a short alternative proof.
\begin{proof} First we  can assume that the class $\{\theta\}$ is big. Indeed, if this is not the case we can just replace $\theta$ with $\theta+\varepsilon\omega$. Using the multi-linearity
of the non-pluripolar product (\cite[Proposition 1.4]{boucksom2010monge}), we can let $\varepsilon\to 0$ to conclude. 

We may assume that $0\geq \varphi,\psi$. 
We divide the proof into two steps.

\noindent {\bf Step 1:} We assume that $\theta$ is semi-positive.

We argue the same as in the proof of \cite[Theorem 4.1]{Ahag09-MA-pluripolar}. Fix $\varepsilon>0$ sufficiently small.
We set
$$\varphi_k = \max (\varphi, - k - 1 ),\quad u_k = (1 - 2\varepsilon ) \varphi_k - 2\varepsilon k,\; \text{and}\quad v_k = \max ( u_k , \psi - \varepsilon k ).$$
We set $E_k := \{ \varphi < - k \} \cap \{ \varphi > \psi -\varepsilon k \}$. We observe that  $\varphi_k+k<0$ on $E_k$, it follows that
$$(1 - 2\varepsilon ) \varphi_k - 2\varepsilon k=\varphi_k-2\varepsilon(\varphi_k+k)>\psi-\varepsilon k $$ there, so $v_k=(1 - 2\varepsilon ) \varphi_k - 2\varepsilon k$ on $E_k$. By plurifine locality (Lemma \ref{lem: plurifine}), we have
$$\theta_{v_k}^n = (\theta + (1-2\varepsilon )dd^c \varphi_k )^n = ((1 -  2\varepsilon ) \theta_{ \varphi_k } + 2\varepsilon\theta )^n$$
on  $E_k$. It follows that
%$$\int_{ E_k } \theta_{ v_k }^n \geq ( 1 - 2\varepsilon )^n \int_{ E_k }  \theta_{ \varphi_k }^n.$$
%Hence, we have
$$\int_{ \{u_k> \psi - \varepsilon k  \} } \theta_{ v_k }^n  \geq ( 1 - 2\varepsilon )^n \int_{ E_k }  \theta_{ \varphi_k }^n.$$
On the other hand, Stoke's theorem yields
$\int_{ X } \theta_{ v_k }^n = \int_{ X } \theta_{ \varphi_k }^n = \int_{ X } \theta^n.$ Observe that $v_k=\psi-\varepsilon k$ on $\{u_k\leq \psi-\varepsilon k \}$.
Therefore, by plurifine locality, we obtain
\begin{equation}\label{eq: thm31}
	\begin{split}
	(1 - 2\varepsilon )^{ n } \int_{ X\backslash E_k }  \theta_{ \varphi_k } ^n   + O(1) \epsilon \int_{ X } \theta^n  \geq \int_{ \{ u_k \leq  \psi - \varepsilon k \} }  \theta_{ \psi }^n  \geq \int_{ \{ \psi \geq  - \varepsilon k \} }  \theta_{ \psi }^n 
	\end{split}
\end{equation} noticing that $\varphi_k\leq 0$. By plurifine locality again, we have
\[ \int_{ X\backslash E_k }  \theta_{ \varphi_k } ^n \leq \int_{\{\varphi>-k-1\}}\theta_\varphi^n+\int_{\{\varphi\leq \psi-\varepsilon k\}}\theta_{\varphi_k}^n.\]
Plugging this into~\eqref{eq: thm31}, we obtain
\[ \int_{\{\varphi>-k-1\}} \theta_\varphi^n+\int_{\{\varphi\leq \psi -\varepsilon k\}}\theta_{\varphi_k}^n+O(1)\varepsilon\int_X\theta^n\geq \int_{\{\psi\geq -\varepsilon k\}}\theta_\psi^n\]
 Since $\varphi\geq \psi+ O(1)$, letting $k\to+\infty$ and then $\varepsilon\to 0^+$, we  obtain 
$$\int_{ X }  \theta_{ \varphi }^n  \geq \int_{ X } \theta_{ \psi }^n$$ since the measures $\theta_\varphi^n$ and $\theta_\psi^n$ vanish on pluripolar sets.
The desired equality follows by reversing the roles of $\varphi$ and $\psi$.

\smallskip
\noindent {\bf Step 2:} We treat the general case: $\{\theta \} $ is merely big.

%and $$\lim_{ k \to +\infty }k^n \capK_{ \omega } ( \{ | \varphi - \psi | > -ch + k \} ) = 0,$$ with $h\in \mathcal E(X,\omega )$. 

We choose a real number $t_0>0$ so large that $\theta + t_0 \omega \geq 0$. By step 1, we have
$$\int_{ X }  (t\omega + \theta_{ \varphi } )^n  = \int_{ X } (t\omega + \theta_{ \psi } )^n ,$$
for all $t\geq t_0$. 
Since non-pluripolar products are multilinear; see \cite[Proposition 1.4]{boucksom2010monge}), we see that $\int_{ X }  (t\omega + \theta_{ \varphi } )^n  \;\text{and}\; \int_{X}  (t\omega + \theta_{ \psi } )^n $ are both homogeneous polynomials of degree $n$. Therefore, we obtain an equality between two polynomials for all $t\geq 0$. Identifying the constant coefficients of these polynomials yields the desired equality.
 \end{proof}
 Next, we establish a slight generalization of Witt-Nystr\"om's result, which is used to prove Theorem~\ref{mainthm}.
\begin{theorem}\label{thm: sub-main} Let $\varphi,\psi\in\PSH(X,\theta)$. If $\varphi$ and $\psi$ have the same singularity type in capacity, i.e.,
    there exist $0\geq h\in\mathcal{E}(X,\omega)$ and $c\geq 0$ such that 
	\[\lim_{k\to+\infty}k^n\capK_{\omega}(\{|\varphi-\psi|>-ch+k\})=0,  \]
	then $\int_X\theta_\varphi^n=\int_X\theta_\psi^n$.
\end{theorem}
\begin{proof} According to Proposition~\ref{prop1}, it suffices to treat the case where $\theta$ is semi-positive.
% First we  can assume that the class $\{\theta\}$ is big. Indeed, if this is not the case we can just replace $\theta$ with $\theta+\varepsilon\omega$. Using the multi-linearity
% of the non-pluripolar product (\cite[Proposition 1.4]{boucksom2010monge}), we can let $\varepsilon\to 0$ to conclude.   
% We may assume that $0\geq \varphi,\psi$. 
% We divide the proof into two steps.
%     \smallskip
% \noindent {\bf Step 1:} We assume that $\theta$ is semi-positive. 
 %We will give a direct proof of the theorem in the case of $\varphi_1 = ... = \varphi_n =\varphi$, $\psi_1 = ... = \psi_n = \psi$ and $\theta_1 = ... =\theta_k = \theta$ is a semi-positive form. 
% We argue the same as in the proof of \cite[Theorem 4.1]{Ahag09-MA-pluripolar}. Fix $\varepsilon>0$ sufficiently small.
% We set
% $$\varphi_k = \max (\varphi, - k - 1 ),\quad u_k = (1 - 2\varepsilon ) \varphi_k - 2\varepsilon k,$$
% and
% $$v_k = \max ( u_k , \psi - \varepsilon k ).$$
% We set $E_k := \{ \varphi < - k \} \cap \{ \varphi > \psi -\varepsilon k \}$. We observe that  $\varphi_k+k<0$ on $E_k$, it follows that
% $$v_k = (1 - 2\varepsilon ) \varphi_k - 2\varepsilon k$$there. By plurifine locality (Lemma \ref{lem: plurifine}), we have
% $$\theta_{v_k}^n = (\theta + (1-2\varepsilon )dd^c \varphi_k )^n = ((1 -  2\varepsilon ) \theta_{ \varphi_k } + 2\varepsilon\theta )^n$$
% on the set $E_k$. It follows that
% %$$\int_{ E_k } \theta_{ v_k }^n \geq ( 1 - 2\varepsilon )^n \int_{ E_k }  \theta_{ \varphi_k }^n.$$
% %Hence, we have
% $$\int_{ \{u_k> \psi - \varepsilon k  \} } \theta_{ v_k }^n  \geq ( 1 - 2\varepsilon )^n \int_{ E_k }  \theta_{ \varphi_k }^n.$$
% On the other hand, Stoke's theorem yields
% $$\int_{ X } \theta_{ v_k }^n = \int_{ X } \theta_{ \varphi_k }^n = \int_{ X } \theta^n.$$ Observe that $v_k=\psi-\varepsilon k$ on $\{u_k\leq \psi-\varepsilon k \}$.
% Therefore, by plurifine locality, we obtain
We borrow the notation in the proof of Proposition~\ref {prop1}. By the same arguments (see~\eqref{eq: thm31}), we have
\begin{equation}\label{eq}
	\begin{split}
	(1 - 2\varepsilon )^{ n } \int_{ X\backslash E_k }  \theta_{ \varphi_k } ^n   + O(1) \epsilon \int_{ X } \theta^n \geq \int_{ \{ \psi \geq  - \varepsilon k \} }  \theta_{ \psi }^n.
	\end{split}
 \end{equation} 
 %noticing that $\varphi_k\leq 0$.
Since $\theta\leq \omega$, we have
\begin{equation*}\label{eq: eq1}
\begin{split}
\int_{ X\backslash E_k } \theta_{ \varphi_k }^n&\leq\int_{ \big\{\varphi > - k - 1 \big\} }  \theta_{ \varphi } ^n + \int_{ \big\{ \varphi \leq \psi - \varepsilon k \big\} } \theta_{ \varphi_k }^n\\
&\leq \int_{X}\theta_{ \varphi }^n  + \int_{ \{ \varphi < \psi + ch - \frac { \varepsilon ( k - 1 ) } { 2 } \} } \theta_{ \varphi_k }^n + \int_{ \big\{ h \leq - \frac { \varepsilon k } { 2c } \big\} } \theta_{ \varphi_k }^n\\
&\leq \int_{ X }  \theta_{ \varphi }^n  +  ( k + 1 )^n \capK_{ \omega } \bigg( \bigg\{ \varphi < \psi + ch - \frac { \varepsilon ( k - 1 ) } { 2 } \bigg\} \bigg) + \int_{ \{ h \leq - \frac { \varepsilon k } { 2c } \} }  \theta_{ \varphi_k }^n.
	\end{split}
\end{equation*}
The last term is dominated by
\begin{equation*}\label{eq: eq2}
\begin{split}
\int_{ \big\{ h\leq - \frac {\varepsilon k } { 2c } \big\} } \theta_{ \varphi_k }^n & \leq \int_{ \big\{ h < - \frac { \varepsilon k } { 4c } + \frac { \varepsilon ( k - 1 ) } { 4c ( k + 1 ) } \varphi_k \big\} } \theta_{ \varphi_k }^n\\
  &\leq \frac { (4c ( k + 1 ) )^n } {  (k - 1 )^n \varepsilon^n } \int_{ \big\{ h < - \frac { \varepsilon k } { 4c } + \frac { \varepsilon ( k - 1 ) } { 4c ( k + 1 ) } \varphi_k\big \} } \bigg( \theta+\ddc {  \frac { \varepsilon ( k - 1 ) } { 4c ( k + 1 ) } \varphi_k }  \bigg)^n\\
&\leq \frac { ( 4c ( k + 1 ) )^n } { ( k - 1 )^n \varepsilon^n } \int_{ \big\{ h < - \frac { \varepsilon k } { 4c } + \frac { \varepsilon ( k - 1 ) } { 4c ( k + 1 ) } \varphi_k \big\} } \bigg( \omega+\ddc {  \frac { \varepsilon ( k - 1 ) } { 4c ( k + 1 ) } \varphi_k }  \bigg)^n\\
&\leq \frac { ( 4c ( k + 1 ) )^n } { ( k - 1 )^n \varepsilon^n } \int_{ \big \{ h < - \frac { \varepsilon k } { 4c } + \frac { \varepsilon ( k - 1 ) } { 4c ( k + 1 ) } \varphi_k \big\} } (\omega + dd^ch )^n\\
&\leq \frac { ( 4c ( k + 1 ) )^n } { ( k - 1 )^n \varepsilon^n } \int_{ \big\{ h < - \frac { \varepsilon k } { 4c } \big\} } (\omega + dd^ch )^n,
\end{split}
\end{equation*} where in the fourth inequality, we applied the comparison principle for the class $\mathcal E(X, \omega)$ (see \cite[Theorem 1.5]{guedj2007weighted} or \cite[Theorem 2.3]{dinew2009uniqueness}).
Combining these inequalities, we obtain
\begin{equation}\label{eq1}
\begin{split}
	\int_{ X\backslash E_k } \theta_{ \varphi_k }^n&\leq \int_{ X }  \theta_{ \varphi }^n  + ( k + 1 )^n \capK_{ \omega } \bigg( \bigg\{ \varphi < \psi + ch - \frac { \varepsilon ( k - 1 ) } { 2 } \bigg\} \bigg) \\
	&\quad +\frac { ( 4c ( k + 1 ) )^n } { ( k - 1 )^n \varepsilon^n } \int_{ \big\{ h < - \frac { \varepsilon k } { 4c } \big\} } (\omega + dd^ch )^n.
\end{split}
\end{equation}
Plugging \eqref{eq1} into \eqref{eq}, we obtain
\begin{align*}
&\int_{ X } \theta_{ \varphi }^n  + 
(k + 1 )^n \capK_{ \omega } \bigg( \bigg\{ \varphi < \psi + ch - \frac { \varepsilon ( k - 1 ) } { 2 } \bigg\} \bigg) 
 \\&\quad  
 +\frac { ( 4c ( k + 1 ) )^n } { ( k - 1 )^n \varepsilon^n } \int_{ \{ h < - \frac { \varepsilon k } { 4c } \} } (\omega + dd^ch )^n   + O(1) \varepsilon \int_{ X } \theta^n\\
&\qquad\geq \int_{ \{ \psi \geq  - \varepsilon k \} } \theta_{ \psi }^n.\end{align*}
% Letting $k\to +\infty$, since the measures $ \theta_{\varphi}^n $, $\theta_{\psi}^n $, and $\omega_h^n$ vanish on pluripolar sets, we obtain
% $$\int_{ X }  \theta_{ \varphi }^n + O(1)\varepsilon \int_{ X } \theta^n\geq \int_{ X } \theta_{ \psi }^n .$$
Since the measures $ \theta_{\varphi}^n $, $\theta_{\psi}^n $, and $\omega_h^n$ vanish on pluripolar sets, letting $k\to +\infty$, and then $\epsilon\to 0^+$ we obtain
$$\int_{ X }  \theta_{ \varphi }^n  \geq \int_{ X } \theta_{ \psi }^n.$$
The desired equality follows by reversing the roles of $\varphi$ and $\psi$. \end{proof}
% Since $\varphi$ and $\psi$ play the same role we get the desired equality.

% \smallskip
% \noindent {\bf Step 2:} We treat the general case: $\{\theta \} $ is merely big.

% %and $$\lim_{ k \to +\infty }k^n \capK_{ \omega } ( \{ | \varphi - \psi | > -ch + k \} ) = 0,$$ with $h\in \mathcal E(X,\omega )$. 

% We choose a real number $t_0>0$ so large that $\theta + t_0 \omega \geq 0$. By step 1, we have
% $$\int_{ X }  (t\omega + \theta_{ \varphi } )^n  = \int_{ X } (t\omega + \theta_{ \psi } )^n ,$$
% for all $t\geq t_0$. 
% Since non-pluripolar products are multilinear; see \cite[Proposition 1.4]{boucksom2010monge}), we see that $$\int_{ X }  (t\omega + \theta_{ \varphi } )^n  \;\text{and}\; \int_{X}  (t\omega + \theta_{ \psi } )^n $$ are both homogeneous polynomials of degree $n$. Therefore, we obtain an equality between two polynomials for all $t\geq 0$. Identifying the constant coefficients of these polynomials yields the desired equality.

\begin{proof}[Proof of Theorem \ref{mainthm}]	
 Arguing the same as in the proof of Proposition~\ref{prop1}, we can assume that the classes $\{\theta_j\}$ are big.

 We first consider the case that $\varphi_j$ have the same singularity type in capacity as $\psi_j$ for $1\leq j\leq n$, i.e.,
$$\lim\limits_{ k \to +\infty }k^n \capK_{ \omega } ( \{ |\varphi_j - \psi_j | > -ch + k \} ) = 0,$$
 for some $c\geq 0$, and $ h\in \mathcal E(X,\omega)$. By Theorem \ref{thm: sub-main}, we have
$$\int_{ X }  \Big(\sum\limits_{ j = 1 }^{ n } t_j \theta_{ j, \varphi_j } \Big)^n  = \int_{ X }  \Big( \sum\limits_{ j = 1 }^{ n } t_j \theta_{ j, \psi_j } \Big)^n ,$$
for all $t_1, \ldots, t_n\geq 0$. Using the 
multilinearity of the non-pluripolar product (see \cite[Proposition 1.4]{boucksom2010monge}), we see that both $\int_{ X }  ( \sum_{ j = 1 }^{ n } t_j \theta_{ \varphi_j } )^n $ and $\int_{ X }  (\sum_{ j = 1 }^{ n } t_j \theta_{ \psi_j } )^n $ are homogeneous polynomials of degree $n$. Therefore, our last identity ensures that all the coefficients of these polynomials have to be equal, which implies
$$\int_X  \theta_{ 1, \varphi_1 }\wedge \ldots \wedge \theta_{ n, \varphi_n }  = \int_X  \theta_{ 1, \psi_1 }\wedge \ldots\wedge \theta_{ n, \psi_n }.$$
We now treat the general case. We set
$$u_{j\ell} = \max ( \varphi_j - \ell, \psi_j ).$$
We see that $\varphi_j$ and $u_{j\ell}$ have the same singularity type in capacity for $j=1,\ldots,n$. The previous case yields
$$\int_X  \theta_{ 1, \varphi_1 }\wedge ...\wedge \theta_{ n, \varphi_n } = \int_X  \theta_{ 1, u_{1\ell} } \wedge ...\wedge \theta_{ n, u_{n\ell} }  .$$
Letting $\ell\to +\infty$, by \cite[Theorem 2.3]{darvas2018monotonicity}
  (see \cite[Proposition 2.7]{Hiep10-CV-capacity} in the local setting), we obtain
\[\int_X  \theta_{1, \varphi_1 }\wedge \ldots  \wedge \theta_{ n, \varphi_n } \geq \int_X \theta_{ 1, \psi_1 }\wedge \ldots\wedge \theta_{ n, \psi_n } .\]\end{proof}

As a direct consequence, we have the following.
\begin{corollary}\label{thm: compare-mass} Let $\{\theta\}$ be a big cohomology class.
Assume $\varphi,\psi\in\PSH(X,\theta)$ such that $\varphi\leq \psi$. If $\int_X\theta_\varphi^n=\int_X\theta_\psi^n$ then $P_\omega(b(\varphi-\psi))\in\mathcal E(X,\omega)$ for all $b>0$.
 
		 Conversely,  if there exists a constant $b>0$ such that $P_\omega(b(\varphi-\psi))\in\mathcal E(X,\omega)$ then $\int_X\theta_\varphi^n=\int_X\theta_\psi^n$.	
\end{corollary}


\begin{proof}
Since $\int_X\theta_\varphi^n=\int_X\theta_\psi^n$, it follows that $\varphi\in\mathcal{E}(X,\theta,\psi)$. Consequently, Proposition~\ref{prop: homega} ensures that $P_\omega(b(\varphi-\psi))\in\mathcal{E}(X,\omega)$ for any $b>0$.

In contrast, setting $h=P_\omega(b(\varphi-\psi))$ for some $b>0$, we observe that
$\varphi\geq \psi +b^{-1}h$ on $X$.  Applying Theorem~\ref{mainthm}, we obtain the desired inequality. 
\end{proof}
% \begin{remark}\label{rmk: mass}
% The second statement can be proved by applying the monotonicity result of Witt-Nystr\"om. Indeed, since $\psi\geq \varphi\geq\psi+ch$ for $c>0$ and $h\in\mathcal{E}(X,\omega)$, the monotonicity of non-pluripolar masses~\cite{wittnystrom19-monotonicity} yields
% \begin{align*}
% \int_X(\theta+c\omega+\ddc\psi)^n&\geq \int_X(\theta+c\omega+\ddc\varphi)^n\\
% & \geq \int_X (\theta+ c\omega+\ddc\psi +c\ddc h)^n\\
% &=\sum_{k=0}^n \binom{n}{k} \theta_\psi^k\wedge c^{n-k}\omega_h^{n-k}\\
% &=\sum_{k=0}^n \binom{n}{k} \theta_\psi^k\wedge c^{n-k}\omega^{n-k}\\
% &=\int_X(\theta+c\omega+\ddc\psi)^n
% \end{align*} where the forth line follows from \cite[Proposition 3.1]{darvas2018monotonicity}. 
% Hence, the inequalities above are, in fact, equalities. It follows that $\int_X\theta_\varphi^k\wedge \omega^{n-k}=\int_X\theta_\psi^k\wedge \omega^{n-k}$, for $k\in\{0,1,\ldots,n \}$, so we are done.
% \end{remark}

% Applying~Theorem \ref{mainthm}, we obtain
% \[\int_X\theta_\varphi^n\geq \int_X\theta_\psi^n .\]
% This yields equality since $\varphi\leq \psi$. 


% We obtain a slight generalization of the first statement of Theorem~\ref{thm: compare-mass} whose proof follows in the same way, using Proposition~\ref{prop: hOmega} instead.
% \begin{theorem}\label{thm: classE}
% Let $\varphi,\psi\in\PSH(X,\theta)$ be such that $P_\theta[\varphi]\simeq P_\theta[\psi]$. Then, for any $b>0$, $P_\omega(b(\varphi-\psi))\in\mathcal E(X,\omega)$ and $P_\omega(b(\psi-\varphi))\in\mathcal E(X,\omega)$.
% \end{theorem}
% By the above result, we obtain the proof of Theorem~\ref{thm: thm1}.
Applying Theorem~\ref{mainthm}, we provide the proof of Theorem~\ref{thm: thm1}.
\begin{proof}[Proof of Theorem~\ref{thm: thm1}]
Set $b\colonequals 2c^{-1}$. Since $\varphi,\psi\in\mathcal{E}(X,\theta,\phi)$, by Proposition~\ref{prop: homega}, both functions $u_b\colonequals P_\omega(b(\varphi-\psi))$ and $v_b\colonequals P_\omega(b(\psi-\varphi))$ belong to the class $\mathcal{E}(X,\omega)$. We set
$h_b\colonequals\frac{u_b+v_b-C_b}{2}$, where $C_b>0$ is an upper bound for $u_b$ and $v_b$. We observe that
\[ch_b=b^{-1}(u_b+v_b-C_b)\leq \varphi-\psi \] since $v_b-C_b\leq 0$.  Similarly, $ch_b\leq \psi-\varphi$, and the first statement follows.   
%Set $h_b=P_\omega(\min(u_b,v_b))$. We observe that $h_b\in\mathcal{E}(X,\omega)$ as follows from~\cite[Proposition 5.3]{darvas2020metric} (see also~\cite[Corollary 3.22]{lu2022hessian}). This completes the proof.

The second statement follows immediately from Theorem~\ref{mainthm}, while the last one is the special case where $\psi=\phi$.  
\end{proof}

\begin{example}
We provide here an example to show that in Theorem~\ref{thm: thm1}, the bigness of $\theta$ is necessary. Let $X=X_1\times X_2$ where $(X_1,\omega_1)$, $(X_2,\omega_2)$ are two compact K\"ahler manifolds of dimension $p\geq 1$, $q\geq 1$ respectively, with $n=\dim X=p+q$.  
Define a K\"ahler form on $X$ by $\omega=\pi_1^*\omega_1+\pi_2^*\omega_2$ is , where $\pi_j$ is the projection onto the $j$-th factor.
The smooth closed (1,1) form $\theta$ is defined by $\pi_1^*\omega_1$. We easily see that $V_\theta=0$ and $\{\theta\}$ is not big since $\Vol(\theta)=\int_X\theta^n=0$. 
% We have the following identity \[\PSH(X_1,\omega_1)\simeq \PSH(X,\theta)|_{X_1}. \]
For any $\varphi\in\PSH(X_1,\omega_1)$, we can find an extension $\Tilde{\varphi}$  of $\varphi$ which belongs to $\PSH(X,\theta)$. 
We have $\int_X\theta_{\Tilde{\varphi}}^n=0$, hence $\Tilde{\varphi}\in\mathcal{E}(X,\theta)$. If Theorem~\ref{thm: thm1} holds, it follows that $\nu(\Tilde\varphi,x)=\nu(V_\theta,x)=0$ for all $x\in X$. This leads to a contradiction since  $\varphi\in\PSH(X_1,\omega_1)$ can be chosen to have a positive Lelong number at some point.
\end{example}
% As an application, the next corollary describes the local/global singular behavior of potentials in the relative full mass class $\mathcal{E}(X,\theta,\phi)$, which provides a slight generalization as well as an alternative proof to the result of Darvas--Di Nezza--Lu~\cite[Theorem 1.1]{darvas2018singularity}.  

% \begin{corollary}\label{thm: ddl18}
% Let $(X, \omega)$ be a compact K\"{a}hler manifold and $\theta$ a smooth closed real (1, 1)-form on $X$ whose cohomology class is pseudoeffective. Let $\phi\in\PSH(X,\theta)$. 
% \begin{enumerate}
% \item If $\varphi,\psi \in\mathcal{E}(X,\theta,\phi)$, then we have
% \[\nu(\varphi,x)=\nu(\psi,x)\quad\text{\and}\quad\mathcal{I}(t\varphi,x)=\mathcal{I}(t\psi,x), \; x\in X, t>0  \]
% where $\nu(\varphi,x)$ is the Lelong number of $\varphi$ at $x$ and $\mathcal{I}(t\varphi,x)$ is the multiplier ideal sheaf of $t\varphi$ at $x$.

% % In particular, if $\{\theta\}$ is big and nef class, then $\nu(\varphi,x)=0$ for any $\varphi\in\mathcal{E}(X,\theta)$, $x\in X$.
% \item 
% If $\{\eta\}$ is a big and nef class then
% \[\mathcal{E}(X,\eta)\cap\PSH(X,\theta,\phi)\subset \mathcal{E}(X,\theta,\phi). \]
% In particular, $\nu(\varphi,x)=0$, for any $\varphi\in\mathcal{E}(X,\eta)$, $x\in X$.
% \end{enumerate}

% \end{corollary}

% \begin{proof}
% % [Proof of Theorem~\ref{thm: ddl18}]
% We prove $(1)$.
% Observe that a function $h\in\mathcal{E}(X,\omega)$ has zero Lelong number at all points, by~\cite[Corollary 1.8]{guedj2007weighted}. Theorem~\ref{thm: thm1} tells us that there exists a function $h\in\mathcal{E}(X,\omega)$ such that 
% \[\varphi\geq \psi+h \quad\text{and}\quad \psi\geq \varphi+h.\]
% This yields $\nu(\varphi,x)=\nu(\psi,x)$ for any $x\in X$. 

% If $\{\theta\}$ is big and nef, then $\nu(V_\theta,x)=0$, for any $x\in X$, as follows from~\cite[Propositions 3.2, 3.4]{boucksom2004divisorial}. Therefore, $\nu(\varphi,x)=0$ for any $x\in X$, $\varphi\in\mathcal{E}(X,\theta)$.

% The equality of multiplier ideal sheaves is a combination of Theorem~~\ref{thm: thm1} and Guan--Zhou's resolution of the Demailly strong openness conjecture~\cite{guan2015proof}. Indeed, for $f\in\mathcal{I}(t,\varphi,x)$, $x\in X$, we have $\int_V |f|^2e^{-t\varphi}\omega^n<\infty$ for some open set $x\in V\subset X$. From~\cite[Theorem 1.1]{guan2015proof}, there is $p>1$ such that $\int_V |f|^2e^{-pt\varphi}\omega^n<\infty$, shrinking $V$ if necessary. Set $q:=\frac{p}{p-1}$.
% Since $\psi\geq \varphi+h$ for some $h\in\mathcal{E}(X,\omega)$, 
% H\"older's inequality yields
% \[\int_V|f|^2e^{-t\psi}\omega^n\leq \bigg(\int_V|f|^2e^{-pt\varphi}\omega \bigg)^{\frac{1}{p}} \bigg(\int_V|f|^2e^{-qth}\omega^n \bigg)^{\frac{1}{q}}<\infty\]
% by Skoda's integrability theorem~\cite{skoda1972sous}. This yields $\mathcal{I}(t\varphi,x)\subset\mathcal{I}(t\psi,x)$.
% Reversing the roles of $\varphi$ and $\psi$, the desired equality follows. 

% For (2), let $\varphi\in \mathcal{E}(X,\eta)\cap \PSH(X,\theta,\phi)$.
% We first assume that $\eta$ is semi-positive. 

%  By Proposition \ref{prop: homega}, we have $P_\omega(\varphi)\in\mathcal{E}(X,\omega)$. Since $\varphi\geq P_\omega(\varphi)$ on $X$, we have
% \[ \varphi\geq \phi-\sup_X\phi +P_\omega(\varphi).\]
% By Theorem \ref{thm: thm1} (or Theorem~\ref{thm: compare-mass}), we obtain $\int_X\theta_\varphi^n=\int_X\theta_\phi^n$. This yields $\varphi\in\mathcal{E}(X,\theta,\phi)$.

% We treat the general case when $\{\eta\}$ is big and nef. We fix $A>1$ so large that $\theta\leq \Tilde{\eta}:=\eta+A\omega$ is K\"ahler. We argue the same as in \cite[page 405]{darvas2018singularity} to show that $V_\eta\in\mathcal{E}(X,\Tilde{\eta})$. By Proposition \ref{prop: homega} again, we have $\varphi\geq V_\eta+ h$ for $h=P_\omega(\varphi-V_\eta)\in\mathcal{E}(X,\omega)$. As follows from Theorem~\ref{thm: thm1}, $\varphi\in\mathcal{E}(X,\Tilde{\eta})$. Since $\Tilde{\eta}\geq 0$, by the previous case, we get $\varphi\in\mathcal{E}(X,\theta,\phi)$.

% In particular, if $\theta=\omega$ and $\phi=0$, then $\varphi\in \mathcal{E}(X,\omega)$. The last statement follows from \cite[Corollary 1.8]{guedj2007weighted}.
% \end{proof}
% Replacing the nefness of $\{\eta\}$ in Theorem \ref{thm: ddl18} (ii) with the assumption $\eta\geq\theta$, we obtain a similar result.


% \begin{corollary}\label{coro: psh}
% Let $\{\eta\}$ and $\{\theta\}$ are big classes such that $\eta\geq \theta$. Let $\phi\in\PSH(X,\theta)$ with $\int_X\theta_\phi^n>0$. Then we have the following
% \begin{enumerate}
% \item $\mathcal{E}(X,\eta)\cap \PSH(X,\theta,\phi)\subset\mathcal{E}
% (X,\theta,\phi)$;
% \item  $\mathcal{E}(X,\eta,\phi)\cap\PSH(X,\theta)\subset \mathcal{E}(X,\theta,\phi)$.
% \end{enumerate}
% \end{corollary}

\begin{corollary}
 Let $\{\eta\}$ and $\{\theta\}$ be big cohomology classes. Let $\phi\in\PSH(X,\theta)$. Then the following are equivalent.
\hfill
\begin{enumerate}
\item $\PSH(X,\eta)\cap\mathcal{E}(X,\theta,\phi)\neq \varnothing$.

\item There exists $u\in\PSH(X,\eta)$, $v\in\mathcal{E}(X,\theta,\phi)$, $c>0$, and $h\in\mathcal{E}(X,\omega)$ such that \[u\geq v+ ch. \]

\item $\mathcal{E}(X,\eta)\cap\PSH(X,\theta,\phi)\subset \mathcal{E}(X,\theta,\phi)$.
\end{enumerate}

\end{corollary}

\begin{proof}
 We see that $(1)\Rightarrow (2)$ is trivial. Indeed, we can take $u=v\in \PSH(X,\eta)\cap\mathcal{E}(X,\theta,\phi)$ and $h=0$.
% Assume $u\in \PSH(X,\eta)\cap\mathcal{E}(X,\theta,\phi)$. Thanks to Proposition~\ref{prop: homega}, we obtain $h\colonequals P_\omega(c^{-1}(u-\phi))\in\mathcal{E}(X,\omega)$. Hence,
% \[u\geq \phi+ch. \]

For $(2)\Rightarrow (3)$, let $\varphi\in \mathcal{E}(X,\eta)\cap \PSH(X,\theta,\phi)$. Thanks to Proposition \ref{prop: homega}, we have $h'\coloneqq P_\omega(\varphi-V_\eta)\in\mathcal{E}(X,\omega)$. By assumption, there exists $u\in\PSH(X,\eta)$, $v\in\mathcal{E}(X,\theta,\phi)$, $c>0$, and $h\in\mathcal{E}(X,\omega)$ such that \[u\geq v+ ch. \]
 It follows that 
 \[\varphi\geq V_\eta+h'\geq u-\sup_X u+h'\geq v+ch+h'-\sup_X u. \]
 Theorem \ref{thm: thm1} yields $\int_X\theta_\varphi^n=\int_X\theta^n_v=\int_X\theta_\phi^n$, hence $\varphi\in\mathcal{E}(X,\theta,\phi)$.

 For $(3)\Rightarrow (1)$, if $\mathcal{E}(X,\eta)\cap\PSH(X,\theta,\phi)=\varnothing$, we are done. Otherwise, $u\in \mathcal{E}(X,\eta)\cap\PSH(X,\theta,\phi)$, so $u\in\mathcal{E}(X,\theta,\phi)$. 
\end{proof}
\begin{remark}
When $\eta\geq \theta$, we obviously have $\PSH(X,\eta)\cap\mathcal{E}(X,\theta,\phi)=\mathcal{E}(X,\theta,\phi)\neq \varnothing$. 
\end{remark}
% \begin{proof}
% (1) Assume $\varphi\in \mathcal{E}(X,\eta)\cap \PSH(X,\theta,\phi)$. By Proposition \ref{prop: homega}, we have $h:=P_\omega(\varphi-V_\eta)\in\mathcal{E}(X,\omega)$. Since $\eta\geq \theta$, we have $V_\eta\geq V_\theta\geq \phi-\sup_X\phi$.
% It follows that 
% \[\varphi\geq V_\eta+h\geq \phi-\sup_X\phi+h. \]
% The conclusion of (1) follows from Theorem \ref{thm: thm1}.

% (2) Assume $\varphi\in \mathcal{E}(X,\eta,\phi)\cap \PSH(X,\theta)$. By Proposition \ref{prop: homega} again, we obtain $h':=P_\omega(\varphi-\phi)\in\mathcal{E}(X,\omega)$. Hence, $\phi\geq \varphi\geq\phi+h'$, and by Theorem \ref{thm: thm1} again, we conclude the corollary.
% \end{proof}
\bibliographystyle{alpha}
\bibliography{bibfile}

\end{document}















